\documentclass[11pt,a4paper]{article}
\usepackage[utf8]{inputenc}
\usepackage[T1]{fontenc}
\usepackage[margin=28mm]{geometry}
\begin{document}

\section*{Gettysburg Address (1863)}

\subsection*{Address}

Four score and seven years ago our fathers brought forth on this continent, a new nation conceived in liberty, and dedicated to the proposition that all men are created equal.

Now we are engaged in a great civil war, testing whether that nation or any nation so conceived and so dedicated, can long endure.
We are met on a great battlefield of that war.
We have come to dedicate a portion of that field as a final resting place for those who here gave their lives, that that nation might live.
It is altogether fitting and proper that we should do this.
But, in a larger sense, we can not dedicate \textasciitilde{} we can not consecrate \textasciitilde{} we can not hallow, this ground.

The brave men living and dead who struggled here, have consecrated it far above our poor power to add or detract.
The world will little note, nor long remember what we say here but it can never forget what they did here.
It is for us the living rather to be dedicated here to the unfinished work which they who fought here have thus far so nobly advanced.

It is rather for us to be here dedicated to the great task remaining before us \textasciitilde{} that from these honored dead we take increased devotion to that cause for which they gave the last full measure of devotion \textasciitilde{} that we here highly resolve that these dead shall not have died in vain \textasciitilde{} that this nation under God shall have a new birth of freedom \textasciitilde{} and that government of the people by the people for the people shall not perish from the earth.

\section*{U.S. Bill of Rights (1791)}

\subsection*{Amendment I}

Congress shall make no law respecting an establishment of religion, or prohibiting the free exercise thereof; or abridging the freedom of speech, or of the press; or the right of the people peaceably to assemble, and to petition the Government for a redress of grievances.

\subsection*{Amendment II}

A well regulated Militia, being necessary to the security of a free State, the right of the people to keep and bear Arms, shall not be infringed.

\subsection*{Amendment III}

No Soldier shall, in time of peace be quartered in any house, without the consent of the Owner, nor in time of war, but in a manner to be prescribed by law.

\subsection*{Amendment IV}

The right of the people to be secure in their persons, houses, papers, and effects, against unreasonable searches and seizures, shall not be violated, and no Warrants shall issue, but upon probable cause, supported by Oath or affirmation, and particularly describing the place to be searched, and the persons or things to be seized.

\subsection*{Amendment V}

No person shall be held to answer for a capital, or otherwise infamous crime, unless on a presentment or indictment of a Grand Jury, except in cases arising in the land or naval forces, or in the Militia, when in actual service in time of War or public danger; nor shall any person be subject for the same offence to be twice put in jeopardy of life or limb; nor shall be compelled in any criminal case to be a witness against himself, nor be deprived of life, liberty, or property, without due process of law; nor shall private property be taken for public use, without just compensation.

\subsection*{Amendment VI}

In all criminal prosecutions, the accused shall enjoy the right to a speedy and public trial, by an impartial jury of the State and district wherein the crime shall have been committed, which district shall have been previously ascertained by law, and to be informed of the nature and cause of the accusation; to be confronted with the witnesses against him; to have compulsory process for obtaining witnesses in his favor, and to have the Assistance of Counsel for his defence.

\subsection*{Amendment VII}

In Suits at common law, where the value in controversy shall exceed twenty dollars, the right of trial by jury shall be preserved, and no fact tried by a jury, shall be otherwise re-examined in any Court of the United States, than according to the rules of the common law.

\subsection*{Amendment VIII}

Excessive bail shall not be required, nor excessive fines imposed, nor cruel and unusual punishments inflicted.

\subsection*{Amendment IX}

The enumeration in the Constitution, of certain rights, shall not be construed to deny or disparage others retained by the people.

\subsection*{Amendment X}

The powers not delegated to the United States by the Constitution, nor prohibited by it to the States, are reserved to the States respectively, or to the people.


\section*{Virginia Declaration of Rights (1776)}

\subsection*{Preamble}

A Declaration of Rights is made by the representatives of the good people of Virginia, assembled in full and free convention which rights do pertain to them and their posterity, as the basis and foundation of government.

\subsection*{Section 1}

That all men are by nature equally free and independent and have certain inherent rights, of which, when they enter into a state of society, they cannot, by any compact, deprive or divest their posterity; namely, the enjoyment of life and liberty, with the means of acquiring and possessing property, and pursuing and obtaining happiness and safety.

\subsection*{Section 2}

That all power is vested in, and consequently derived from, the people; that magistrates are their trustees and servants and at all times amenable to them.

\subsection*{Section 3}

That government is, or ought to be, instituted for the common benefit, protection, and security of the people, nation, or community; of all the various modes and forms of government, that is best which is capable of producing the greatest degree of happiness and safety and is most effectually secured against the danger of maladministration.
And that, when any government shall be found inadequate or contrary to these purposes, a majority of the community has an indubitable, inalienable, and indefeasible right to reform, alter, or abolish it, in such manner as shall be judged most conducive to the public weal.

\subsection*{Section 4}

That no man, or set of men, is entitled to exclusive or separate emoluments or privileges from the community, but in consideration of public services; which, nor being descendible, neither ought the offices of magistrate, legislator, or judge to be hereditary.

\subsection*{Section 5}

That the legislative and executive powers of the state should be separate and distinct from the judiciary; and that the members of the two first may be restrained from oppression, by feeling and participating the burdens of the people, they should, at fixed periods, be reduced to a private station, return into that body from which they were originally taken, and the vacancies be supplied by frequent, certain, and regular elections, in which all, or any part, of the former members, to be again eligible, or ineligible, as the laws shall direct.

\subsection*{Section 6}

That elections of members to serve as representatives of the people, in assembly ought to be free; and that all men, having sufficient evidence of permanent common interest with, and attachment to, the community, have the right of suffrage and cannot be taxed or deprived of their property for public uses without their own consent or that of their representatives so elected, nor bound by any law to which they have not, in like manner, assented, for the public good.

\subsection*{Section 7}

That all power of suspending laws, or the execution of laws, by any authority, without consent of the representatives of the people, is injurious to their rights and ought not to be exercised.

\subsection*{Section 8}

That in all capital or criminal prosecutions a man has a right to demand the cause and nature of his accusation, to be confronted with the accusers and witnesses, to call for evidence in his favor, and to a speedy trial by an impartial jury of twelve men of his vicinage, without whose unanimous consent he cannot be found guilty; nor can he be compelled to give evidence against himself; that no man be deprived of his liberty except by the law of the land or the judgment of his peers.

\subsection*{Section 9}

That excessive bail ought not to be required, nor excessive fines imposed, nor cruel and unusual punishments inflicted.

\subsection*{Section 10}

That general warrants, whereby an officer or messenger may be commanded to search suspected places without evidence of a fact committed, or to seize any person or persons not named, or whose offense is not particularly described and supported by evidence, are grievous and oppressive and ought not to be granted.

\subsection*{Section 11}

That in controversies respecting property, and in suits between man and man, the ancient trial by jury is preferable to any other and ought to be held sacred.

\subsection*{Section 12}

That the freedom of the press is one of the great bulwarks of liberty, and can never be restrained but by despotic governments.

\subsection*{Section 13}

That a well-regulated militia, composed of the body of the people, trained to arms, is the proper, natural, and safe defense of a free state; that standing armies, in time of peace, should be avoided as dangerous to liberty; and that in all cases the military should be under strict subordination to, and governed by, the civil power.

\subsection*{Section 14}

That the people have a right to uniform government; and, therefore, that no government separate from or independent of the government of Virginia ought to be erected or established within the limits thereof.

\subsection*{Section 15}

That no free government, or the blessings of liberty, can be preserved to any people but by a firm adherence to justice, moderation, temperance, frugality, and virtue and by frequent recurrence to fundamental principles.

\subsection*{Section 16}

That religion, or the duty which we owe to our Creator, and the manner of discharging it, can be directed only by reason and conviction, not by force or violence; and therefore all men are equally entitled to the free exercise of religion, according to the dictates of conscience; and that it is the mutual duty of all to practise Christian forbearance, love, and charity toward each other.

\section*{Declaration of Independence (1776)}

\subsection*{Opening}

In Congress, July 4, 1776

The unanimous Declaration of the thirteen united States of America, When in the Course of human events, it becomes necessary for one people to dissolve the political bands which have connected them with another, and to assume among the powers of the earth, the separate and equal station to which the Laws of Nature and of Nature's God entitle them, a decent respect to the opinions of mankind requires that they should declare the causes which impel them to the separation.

We hold these truths to be self-evident, that all men are created equal, that they are endowed by their Creator with certain unalienable Rights, that among these are Life, Liberty and the pursuit of Happiness.--That to secure these rights, Governments are instituted among Men, deriving their just powers from the consent of the governed, --That whenever any Form of Government becomes destructive of these ends, it is the Right of the People to alter or to abolish it, and to institute new Government, laying its foundation on such principles and organizing its powers in such form, as to them shall seem most likely to effect their Safety and Happiness.

Prudence, indeed, will dictate that Governments long established should not be changed for light and transient causes; and accordingly all experience hath shewn, that mankind are more disposed to suffer, while evils are sufferable, than to right themselves by abolishing the forms to which they are accustomed.

But when a long train of abuses and usurpations, pursuing invariably the same Object evinces a design to reduce them under absolute Despotism, it is their right, it is their duty, to throw off such Government, and to provide new Guards for their future security.--Such has been the patient sufferance of these Colonies; and such is now the necessity which constrains them to alter their former Systems of Government.

The history of the present King of Great Britain is a history of repeated injuries and usurpations, all having in direct object the establishment of an absolute Tyranny over these States.
To prove this, let Facts be submitted to a candid world.

\subsection*{Grievances}

He has refused his Assent to Laws, the most wholesome and necessary for the public good.

He has forbidden his Governors to pass Laws of immediate and pressing importance, unless suspended in their operation till his Assent should be obtained; and when so suspended, he has utterly neglected to attend to them.

He has refused to pass other Laws for the accommodation of large districts of people, unless those people would relinquish the right of Representation in the Legislature, a right inestimable to them and formidable to tyrants only.

He has called together legislative bodies at places unusual, uncomfortable, and distant from the depository of their public Records, for the sole purpose of fatiguing them into compliance with his measures.

He has dissolved Representative Houses repeatedly, for opposing with manly firmness his invasions on the rights of the people.

He has refused for a long time, after such dissolutions, to cause others to be elected; whereby the Legislative powers, incapable of Annihilation, have returned to the People at large for their exercise; the State remaining in the mean time exposed to all the dangers of invasion from without, and convulsions within.

He has endeavoured to prevent the population of these States; for that purpose obstructing the Laws for Naturalization of Foreigners; refusing to pass others to encourage their migrations hither, and raising the conditions of new Appropriations of Lands.

He has obstructed the Administration of Justice, by refusing his Assent to Laws for establishing Judiciary powers.

He has made Judges dependent on his Will alone, for the tenure of their offices, and the amount and payment of their salaries.

He has erected a multitude of New Offices, and sent hither swarms of Officers to harrass our people, and eat out their substance.

He has kept among us, in times of peace, Standing Armies without the Consent of our legislatures.

He has affected to render the Military independent of and superior to the Civil power.

He has combined with others to subject us to a jurisdiction foreign to our constitution, and unacknowledged by our laws; giving his Assent to their Acts of pretended Legislation:

For Quartering large bodies of armed troops among us:

For protecting them, by a mock Trial, from punishment for any Murders which they should commit on the Inhabitants of these States:

For cutting off our Trade with all parts of the world:

For imposing Taxes on us without our Consent:

For depriving us in many cases, of the benefits of Trial by Jury:

For transporting us beyond Seas to be tried for pretended offences:

For abolishing the free System of English Laws in a neighbouring Province, establishing therein an Arbitrary government, and enlarging its Boundaries so as to render it at once an example and fit instrument for introducing the same absolute rule into these Colonies:

For taking away our Charters, abolishing our most valuable Laws, and altering fundamentally the Forms of our Governments:

For suspending our own Legislatures, and declaring themselves invested with power to legislate for us in all cases whatsoever.

He has abdicated Government here, by declaring us out of his Protection and waging War against us.

He has plundered our seas, ravaged our Coasts, burnt our towns, and destroyed the lives of our people.

He is at this time transporting large Armies of foreign Mercenaries to compleat the works of death, desolation and tyranny, already begun with circumstances of Cruelty \& perfidy scarcely paralleled in the most barbarous ages, and totally unworthy the Head of a civilized nation.

He has constrained our fellow Citizens taken Captive on the high Seas to bear Arms against their Country, to become the executioners of their friends and Brethren, or to fall themselves by their Hands.

He has excited domestic insurrections amongst us, and has endeavoured to bring on the inhabitants of our frontiers, the merciless Indian Savages, whose known rule of warfare, is an undistinguished destruction of all ages, sexes and conditions.

\subsection*{Conclusion}

In every stage of these Oppressions We have Petitioned for Redress in the most humble terms: Our repeated Petitions have been answered only by repeated injury.
A Prince, whose character is thus marked by every act which may define a Tyrant, is unfit to be the ruler of a free people.

Nor have We been wanting in attentions to our Brittish brethren.
We have warned them from time to time of attempts by their legislature to extend an unwarrantable jurisdiction over us.
We have reminded them of the circumstances of our emigration and settlement here.
We have appealed to their native justice and magnanimity, and we have conjured them by the ties of our common kindred to disavow these usurpations, which, would inevitably interrupt our connections and correspondence.

They too have been deaf to the voice of justice and of consanguinity.
We must, therefore, acquiesce in the necessity, which denounces our Separation, and hold them, as we hold the rest of mankind, Enemies in War, in Peace Friends.

We, therefore, the Representatives of the united States of America, in General Congress, Assembled, appealing to the Supreme Judge of the world for the rectitude of our intentions, do, in the Name, and by Authority of the good People of these Colonies, solemnly publish and declare, That these United Colonies are, and of Right ought to be Free and Independent States; that they are Absolved from all Allegiance to the British Crown, and that all political connection between them and the State of Great Britain, is and ought to be totally dissolved; and that as Free and Independent States, they have full Power to levy War, conclude Peace, contract Alliances, establish Commerce, and to do all other Acts and Things which Independent States may of right do.

And for the support of this Declaration, with a firm reliance on the protection of divine Providence, we mutually pledge to each other our Lives, our Fortunes and our sacred Honor.

\section*{Déclaration des droits de l'homme et du citoyen (1789)}

\subsection*{Préambule}

Les représentants du peuple français, constitués en Assemblée nationale, considérant que l'ignorance, l'oubli ou le mépris des droits de l'homme sont les seules causes des malheurs publics et de la corruption des gouvernements, ont résolu d'exposer, dans une déclaration solennelle, les droits naturels, inaliénables et sacrés de l'homme, afin que cette déclaration, constamment présente à tous les membres du corps social, leur rappelle sans cesse leurs droits et leurs devoirs ; afin que les actes du pouvoir législatif, et ceux du pouvoir exécutif, pouvant être à chaque instant comparés avec le but de toute institution politique, en soient plus respectés ; afin que les réclamations des citoyens, fondées désormais sur des principes simples et incontestables, tournent toujours au maintien de la Constitution et au bonheur de tous.

En conséquence, l'Assemblée nationale reconnaît et déclare, en présence et sous les auspices de l'Être suprême, les droits suivants de l'homme et du citoyen.

\subsection*{Article 1}

Les hommes naissent et demeurent libres et égaux en droits.
Les distinctions sociales ne peuvent être fondées que sur l'utilité commune.

\subsection*{Article 2}

Le but de toute association politique est la conservation des droits naturels et imprescriptibles de l'homme.
Ces droits sont la liberté, la propriété, la sûreté, et la résistance à l'oppression.

\subsection*{Article 3}

Le principe de toute souveraineté réside essentiellement dans la nation.
Nul corps, nul individu ne peut exercer d'autorité qui n'en émane expressément.

\subsection*{Article 4}

La liberté consiste à pouvoir faire tout ce qui ne nuit pas à autrui : ainsi, l'exercice des droits naturels de chaque homme n'a de bornes que celles qui assurent aux autres membres de la société la jouissance de ces mêmes droits.
Ces bornes ne peuvent être déterminées que par la loi.

\subsection*{Article 5}

La loi n'a le droit de défendre que les actions nuisibles à la société.
Tout ce qui n'est pas défendu par la loi ne peut être empêché, et nul ne peut être contraint à faire ce qu'elle n'ordonne pas.

\subsection*{Article 6}

La loi est l'expression de la volonté générale.
Tous les citoyens ont droit de concourir personnellement, ou par leurs représentants, à sa formation.
Elle doit être la même pour tous, soit qu'elle protège, soit qu'elle punisse.
Tous les citoyens étant égaux à ses yeux sont également admissibles à toutes dignités, places et emplois publics, selon leur capacité, et sans autre distinction que celle de leurs vertus et de leurs talents.

\subsection*{Article 7}

Nul homme ne peut être accusé, arrêté ni détenu que dans les cas déterminés par la loi, et selon les formes qu'elle a prescrites.
Ceux qui sollicitent, expédient, exécutent ou font exécuter des ordres arbitraires, doivent être punis ; mais tout citoyen appelé ou saisi en vertu de la loi doit obéir à l'instant : il se rend coupable par la résistance.

\subsection*{Article 8}

La loi ne doit établir que des peines strictement et évidemment nécessaires, et nul ne peut être puni qu'en vertu d'une loi établie et promulguée antérieurement au délit, et légalement appliquée.

\subsection*{Article 9}

Tout homme étant présumé innocent jusqu'à ce qu'il ait été déclaré coupable, s'il est jugé indispensable de l'arrêter, toute rigueur qui ne serait pas nécessaire pour s'assurer de sa personne doit être sévèrement réprimée par la loi.

\subsection*{Article 10}

Nul ne doit être inquiété pour ses opinions, même religieuses, pourvu que leur manifestation ne trouble pas l'ordre public établi par la loi.

\subsection*{Article 11}

La libre communication des pensées et des opinions est un des droits les plus précieux de l'homme : tout citoyen peut donc parler, écrire, imprimer librement, sauf à répondre de l'abus de cette liberté dans les cas déterminés par la loi.

\subsection*{Article 12}

La garantie des droits de l'homme et du citoyen nécessite une force publique : cette force est donc instituée pour l'avantage de tous, et non pour l'utilité particulière de ceux auxquels elle est confiée.

\subsection*{Article 13}

Pour l'entretien de la force publique, et pour les dépenses d'administration, une contribution commune est indispensable : elle doit être également répartie entre tous les citoyens, en raison de leurs facultés.

\subsection*{Article 14}

Tous les citoyens ont le droit de constater, par eux-mêmes ou par leurs représentants, la nécessité de la contribution publique, de la consentir librement, d'en suivre l'emploi, et d'en déterminer la quotité, l'assiette, le recouvrement et la durée.

\subsection*{Article 15}

La société a le droit de demander compte à tout agent public de son administration.

\subsection*{Article 16}

Toute société dans laquelle la garantie des droits n'est pas assurée, ni la séparation des pouvoirs déterminée, n'a point de constitution.

\subsection*{Article 17}

La propriété étant un droit inviolable et sacré, nul ne peut en être privé, si ce n'est lorsque la nécessité publique, légalement constatée, l'exige évidemment, et sous la condition d'une juste et préalable indemnité.

\end{document}
